\addcontentsline{toc}{section}{Заключение}

В ходе проведённого исследования был осуществлён сравнительный анализ современных технологий разработки интеллектуальных систем: языков запросов к данным и инструментов построения и развертывания интеллектуальных решений.

В первом разделе работы были рассмотрены языки декларативных запросов LINQ и SPARQL. Проанализированы их синтаксические особенности, основные конструкции и сферы применения. Установлено, что LINQ ориентирован на объектно-ориентированную экосистему .NET и предоставляет универсальный подход к работе с различными типами источников данных. В то же время SPARQL предназначен для работы с RDF-графами и онтологиями в рамках Semantic Web и обеспечивает мощные средства для формулирования запросов к распределённым семантическим данным. Результаты анализа были формализованы с помощью SCn-кода.

Во втором разделе были исследованы инструменты Microsoft Azure AI и Neo4j. Платформа Azure AI предоставляет облачную инфраструктуру и широкий набор сервисов для построения, обучения и развёртывания моделей искусственного интеллекта, включая языковые модели, когнитивные API и инструменты автоматизации ML. В свою очередь, Neo4j является графовой базой данных, специально предназначенной для хранения, визуализации и анализа семантических связей и графов знаний. Оба инструмента имеют ярко выраженные преимущества и используются для различных классов задач. Информация о каждом инструменте также была формализована на SCn-коде.

Полученные результаты могут быть использованы при проектировании гибридных интеллектуальных систем, сочетающих возможности облачных вычислений и графового анализа, а также при выборе технологий для решения задач анализа данных, построения онтологий и обработки знаний в интеллектуальных приложениях.

